%*************************************
% บทที่ 2
%*************************************
\chapterTitle{2}{แนวคิด ทฤษฎี และงานวิจัยที่เกี่ยวข้อง}

%*******************************************
% แนวคิด ทฤษฎีต่าง ๆ
%*******************************************
\section{แนวคิดพื้นฐานเกี่ยวกับโครงงาน}

% 2.1.1
\subsection{ความสำคัญของการเรียนการสอนด้วยระบบปฏิบัติการลีนุกซ์ในการศึกษาด้านเทคโนโลยีสารสนเทศ}

\hspace*{1.3em} %ย่อหน้า
การเรียนการสอนในสาขาเทคโนโลยีสารสนเทศ โดยเฉพาะในระดับอุดมศึกษา มีความจำเป็นอย่างยิ่งที่จะต้องให้นักศึกษาได้ฝึกฝนการใช้งานระบบปฏิบัติการลีนุกซ์ เนื่องจากระบบปฏิบัติการดังกล่าวเป็นพื้นฐานสำคัญในการทำงานด้านเซิร์ฟเวอร์ การพัฒนาระบบ การจัดการโครงสร้างพื้นฐานด้านไอที และเทคโนโลยีคลาวด์ที่กำลังเป็นที่ต้องการของตลาดแรงงานในปัจจุบัน \parencite{anderson2008theory}

\hspace*{1.3em} %ย่อหน้า
การเรียนรู้และฝึกฝนการใช้งานระบบลีนุกซ์มีความสำคัญในหลายประการ ได้แก่
\begin{mycustomitem}
    \item \textbf{พัฒนาทักษะการจัดการระบบ} ช่วยให้นักศึกษาเข้าใจและสามารถจัดการระบบปฏิบัติการที่ใช้กันอย่างแพร่หลายในโลกธุรกิจ โดยเฉพาะในด้านเซิร์ฟเวอร์และระบบคลาวด์
    \item \textbf{เสริมสร้างความเข้าใจในการทำงานของระบบคอมพิวเตอร์} การใช้งานลีนุกซ์ผ่าน Command Line Interface ช่วยให้นักศึกษาเข้าใจการทำงานของระบบในระดับที่ลึกขึ้น
    \item \textbf{เตรียมความพร้อมสู่การทำงาน} ระบบลีนุกซ์เป็นมาตรฐานในอุตสาหกรรมเทคโนโลยี การฝึกฝนการใช้งานจะช่วยให้นักศึกษามีความพร้อมในการทำงานจริง
    \item \textbf{สนับสนุนการเรียนรู้เทคโนโลยีใหม่} หลายเทคโนโลยีสมัยใหม่ เช่น Docker, Kubernetes, และระบบ DevOps ต่าง ๆ ต้องอาศัยความรู้พื้นฐานเกี่ยวกับลีนุกซ์ \parencite{lee2021container}
\end{mycustomitem}

\hspace*{1.3em} %ย่อหน้า
อย่างไรก็ตาม การจัดเตรียมสภาพแวดล้อมสำหรับการเรียนการสอนด้วยระบบลีนุกซ์นั้นมักเผชิญกับความท้าทายหลายประการ ได้แก่
\begin{mycustomitem}
    \item \textbf{ความยุ่งยากในการติดตั้งและตั้งค่า} การติดตั้งระบบลีนุกซ์บนเครื่องคอมพิวเตอร์ส่วนบุคคลของนักศึกษามักต้องใช้เวลานานและมีความซับซ้อน โดยเฉพาะสำหรับผู้ที่ไม่มีประสบการณ์
    \item \textbf{ข้อจำกัดด้านฮาร์ดแวร์} เครื่องคอมพิวเตอร์ส่วนบุคคลของนักศึกษาอาจมีข้อจำกัดด้านหน่วยความจำหรือพื้นที่จัดเก็บข้อมูล ทำให้การใช้งานเครื่องเสมือนไม่เสถียร
    \item \textbf{ความไม่สม่ำเสมอของสภาพแวดล้อม} นักศึกษาแต่ละคนอาจมีการตั้งค่าและเวอร์ชันของซอฟต์แวร์ที่แตกต่างกัน ทำให้เกิดปัญหาในการเรียนการสอนแบบเดียวกัน
    \item \textbf{การบำรุงรักษาและแก้ไขปัญหา} เมื่อเกิดปัญหาทางเทคนิค อาจารย์และนักศึกษาต้องใช้เวลามากในการแก้ไขปัญหาแทนที่จะมุ่งเน้นไปที่การเรียนรู้เนื้อหาหลัก
    \item \textbf{ค่าใช้จ่ายในการใช้บริการภายนอก} การใช้บริการคลาวด์จากผู้ให้บริการภายนอกมักมีค่าใช้จ่ายที่สูง โดยเฉพาะเมื่อต้องใช้งานตลอดภาคการศึกษา
\end{mycustomitem}



% 2.1.2
\subsection{บทบาทของเทคโนโลยี Virtualization และ Cloud Computing ในการศึกษา}

\hspace*{1.3em} %ย่อหน้า
ในยุคที่เทคโนโลยีสารสนเทศพัฒนาอย่างรวดเร็ว การนำเทคโนโลยี Virtualization และ Cloud Computing มาประยุกต์ใช้ในการศึกษาได้กลายเป็นแนวทางที่สำคัญในการแก้ไขปัญหาและเพิ่มประสิทธิภาพการเรียนการสอน โดยเฉพาะอย่างยิ่งในสาขาเทคโนโลยีสารสนเทศที่ต้องการสภาพแวดล้อมที่หลากหลายและซับซ้อนสำหรับการฝึกปฏิบัติ \parencite{clark2016blended}

\hspace*{1.3em} %ย่อหน้า
เทคโนโลยี Virtualization และ Cloud Computing มีบทบาทสำคัญในการเพิ่มประสิทธิภาพการเรียนการสอนหลายประการ ได้แก่
\begin{mycustomitem}
    \item \textbf{ลดเวลาและความซับซ้อนในการตั้งค่าระบบ} เทคโนโลยีเครื่องเสมือนสามารถสร้างสภาพแวดล้อมที่พร้อมใช้งานได้อย่างรวดเร็ว โดยไม่ต้องให้นักศึกษาแต่ละคนติดตั้งและตั้งค่าระบบด้วยตนเอง
    \item \textbf{เพิ่มความยืดหยุ่นในการจัดการทรัพยากร} สามารถปรับเปลี่ยนขนาดและสเปกของเครื่องเสมือนตามความต้องการของแต่ละรายวิชาหรือโครงงาน โดยไม่ต้องลงทุนซื้อฮาร์ดแวร์เพิ่มเติม
    \item \textbf{รองรับการเรียนรู้แบบ Remote Learning} นักศึกษาสามารถเข้าถึงและใช้งานระบบได้จากที่ไหนก็ได้ ตราบใดที่มีการเชื่อมต่ออินเทอร์เน็ต ทำให้การเรียนการสอนไม่ถูกจำกัดด้วยสถานที่และเวลา
    \item \textbf{สร้างความสม่ำเสมอของสภาพแวดล้อม} ทุกคนจะได้ใช้งานระบบที่มีการตั้งค่าเหมือนกัน ลดปัญหาที่เกิดจากความแตกต่างของสภาพแวดล้อมในเครื่องของแต่ละคน
    \item \textbf{ประหยัดค่าใช้จ่าย} การใช้ทรัพยากรแบบรวมศูนย์ช่วยลดค่าใช้จ่ายในการซื้อและบำรุงรักษาฮาร์ดแวร์ รวมถึงลดค่าใช้จ่ายสำหรับนักศึกษาในการจัดหาเครื่องมือสำหรับการเรียน
\end{mycustomitem}

\hspace*{1.3em} %ย่อหน้า
แนวคิดของแพลตฟอร์มคลัสเตอร์สำหรับการศึกษาจึงเป็นการต่อยอดจากข้อดีของเทคโนโลยีเหล่านี้ โดยการสร้างระบบที่สามารถ
\begin{mycustomitem}
    \item \textbf{จัดการเครื่องเสมือนแบบรวมศูนย์} ผ่านระบบ Web Application ที่ใช้งานง่าย ทำให้อาจารย์และนักศึกษาสามารถร้องขอ สร้าง และจัดการเครื่องเสมือนได้โดยไม่ต้องมีความรู้ลึกด้านระบบ
    \item \textbf{ใช้ประโยชน์จากทรัพยากรที่มีอยู่อย่างเต็มที่} การรวมเครื่องเซิร์ฟเวอร์หลายเครื่องเข้าเป็นคลัสเตอร์เดียว ช่วยให้สามารถใช้งานทรัพยากรได้อย่างมีประสิทธิภาพสูงสุด
    \item \textbf{รองรับการใช้งานหลากหลายรูปแบบ} ทั้งการเรียนการสอนปกติ การทำโครงงาน การวิจัย และการพัฒนาระบบต่าง ๆ ภายในภาควิชา
    \item \textbf{มีระบบติดตามและบริหารจัดการ} ด้วยระบบ Monitoring ที่ช่วยในการติดตามการใช้งานและแก้ไขปัญหาได้อย่างรวดเร็ว
\end{mycustomitem}



%2.2
\section{แนวคิดและทฤษฎีที่เกี่ยวข้องกับแพลตฟอร์มคลัสเตอร์เสมือน}

%2.2.1
\subsection{เทคโนโลยี Virtualization}

\hspace*{1.3em} %ย่อหน้า
Virtualization Technology เป็นเทคนิคการจำลองซอฟต์แวร์หรือฮาร์ดแวร์บนซอฟต์แวร์อื่น โดยสภาพแวดล้อมจำลองนี้เรียกว่าเครื่องเสมือน (Virtual Machine) ซึ่งเป็นเทคโนโลยีพื้นฐานสำคัญที่ทำให้สามารถแบ่งทรัพยากรฮาร์ดแวร์เครื่องเดียวออกเป็นหลาย ๆ ระบบที่ทำงานได้อิสระจากกัน \parencite{garcia2023hypervisor}

\hspace*{1.3em} %ย่อหน้า
\textbf{Hypervisor และประเภทของ Virtualization}
\begin{mycustomitem}
    \item \textbf{Type 1 Hypervisor (Bare-metal)} เป็น Hypervisor ที่ทำงานโดยตรงบนฮาร์ดแวร์ โดยไม่ต้องผ่านระบบปฏิบัติการหลัก เช่น VMware vSphere, Microsoft Hyper-V, และ Proxmox VE ซึ่งมีประสิทธิภาพสูงและเหมาะสำหรับการใช้งานในระดับองค์กร
    \item \textbf{Type 2 Hypervisor (Hosted)} เป็น Hypervisor ที่ทำงานบนระบบปฏิบัติการหลัก เช่น VMware Workstation, VirtualBox ซึ่งง่ายต่อการใช้งานแต่มีประสิทธิภาพต่ำกว่า Type 1
    \item \textbf{Container-based Virtualization} เป็นรูปแบบของ Virtualization ที่แบ่งปันระบบปฏิบัติการเดียวกัน เช่น Docker, LXC ซึ่งมีความเร็วและประสิทธิภาพสูงกว่าแต่มีความแยกตัวน้อยกว่าเครื่องเสมือนแบบเต็มรูปแบบ
\end{mycustomitem}

\hspace*{1.3em} %ย่อหน้า
\textbf{ข้อดีของเทคโนโลยี Virtualization}
\begin{mycustomitem}
    \item \textbf{การใช้ทรัพยากรอย่างมีประสิทธิภาพ} สามารถใช้ทรัพยากรฮาร์ดแวร์ได้เต็มประสิทธิภาพ เนื่องจากสามารถรันหลายระบบพร้อมกันได้บนเครื่องเดียว
    \item \textbf{ความยืดหยุ่นในการจัดการ} สามารถสร้าง ลบ หรือปรับแต่งเครื่องเสมือนได้อย่างรวดเร็ว โดยไม่ต้องแตะต้องฮาร์ดแวร์จริง
    \item \textbf{การแยกตัวและความปลอดภัย} แต่ละเครื่องเสมือนทำงานแยกจากกัน หากเครื่องหนึ่งเกิดปัญหาจะไม่ส่งผลกระทบต่อเครื่องอื่น
    \item \textbf{ความสะดวกในการบำรุงรักษา} สามารถสำรองข้อมูล ย้าย หรือกู้คืนระบบได้ง่าย โดยไม่ต้องหยุดการทำงานของเครื่องอื่น
\end{mycustomitem}


%2.2.2
\subsection{Cloud Computing และ Container Orchestration}

\hspace*{1.3em} %ย่อหน้า
Cloud Computing เป็นแนวคิดการประมวลผลแบบคลาวด์ที่ช่วยให้สามารถเข้าถึงทรัพยากรการประมวลผลที่กำหนดค่าได้ตามต้องการ ซึ่งสามารถจัดเตรียมและปล่อยทรัพยากรได้อย่างรวดเร็วโดยใช้ความพยายามในการจัดการหรือการโต้ตอบระหว่างผู้ให้บริการน้อยที่สุด โดยมีรูปแบบการให้บริการหลัก ได้แก่

\hspace*{1.3em} %ย่อหน้า
\textbf{รูปแบบการให้บริการ Cloud Computing}
\begin{mycustomitem}
    \item \textbf{Infrastructure as a Service (IaaS)} การให้บริการโครงสร้างพื้นฐานด้านไอที เช่น เครื่องเสมือน พื้นที่จัดเก็บข้อมูล เครือข่าย โดยผู้ใช้สามารถควบคุมระบบปฏิบัติการและแอปพลิเคชันได้
    \item \textbf{Platform as a Service (PaaS)} การให้บริการแพลตฟอร์มสำหรับการพัฒนาและปรับใช้แอปพลิเคชัน โดยไม่ต้องจัดการโครงสร้างพื้นฐานด้านล่าง
    \item \textbf{Software as a Service (SaaS)} การให้บริการซอฟต์แวร์ที่พร้อมใช้งานผ่านอินเทอร์เน็ต โดยไม่ต้องติดตั้งหรือจัดการด้วยตนเอง
\end{mycustomitem}

\hspace*{1.3em} %ย่อหน้า
\textbf{Container และ Container Orchestration}

\hspace*{1.3em} %ย่อหน้า
Container คือการรวมแอปพลิเคชันและรหัสไบนารี ไลบรารี และการกำหนดค่าเข้าด้วยกันในขณะที่แบ่งปันอิมเมจระบบปฏิบัติการโฮสต์ ดังนั้นคอนเทนเนอร์จึงแบ่งปันทรัพยากรและดำเนินการไมโครเซอร์วิสขนาดเล็ก โปรแกรมซอฟต์แวร์ และแอปพลิเคชันที่ครอบคลุมยิ่งขึ้นได้อย่างมีประสิทธิภาพด้วยค่าใช้จ่ายด้านการจัดการน้อยกว่าเครื่องเสมือน

\begin{mycustomitem}
    \item \textbf{ข้อดีของ Container} มีขนาดเล็ก เริ่มต้นเร็ว ใช้ทรัพยากรน้อย และสามารถพกพาได้ง่าย เนื่องจากบรรจุทุกสิ่งที่จำเป็นสำหรับการทำงานของแอปพลิเคชัน
    \item \textbf{Docker} เป็นแพลตฟอร์ม Container ที่ได้รับความนิยมสูงสุด ช่วยให้การสร้าง จัดการ และปรับใช้ Container เป็นเรื่องง่าย
\end{mycustomitem}

\hspace*{1.3em} %ย่อหน้า
Container Orchestration คือการปรับใช้และจัดการแอปพลิเคชันที่เป็น Container ที่อยู่ในคลัสเตอร์ที่เชื่อมต่อถึงกัน โดยมีโหนดที่คอยควบคุมการทำงานที่เรียกว่าโหนดมาสเตอร์ (Master node) ซึ่งมีหน้าที่ดูแลคลัสเตอร์และจัดการการปรับใช้คอนเทนเนอร์บนโหนดเวิร์กเกอร์ (Worker node)

\begin{mycustomitem}
    \item \textbf{การจัดการทรัพยากรอัตโนมัติ} ระบบจะจัดสรรทรัพยากรให้กับ Container ตามความต้องการและสถานะของคลัสเตอร์
    \item \textbf{High Availability} สามารถกระจาย Container ไปยังหลายโหนดเพื่อป้องกันการหยุดทำงานเมื่อโหนดใดโหนดหนึ่งเกิดปัญหา
    \item \textbf{Auto Scaling} สามารถปรับขนาดของแอปพลิเคชันขึ้นลงอัตโนมัติตามปริมาณการใช้งาน
    \item \textbf{Service Discovery และ Load Balancing} ช่วยในการเชื่อมต่อ Container ต่าง ๆ และกระจายโหลดการทำงาน
\end{mycustomitem}
%2.2.3
\subsection{Web Application และเทคโนโลยีสนับสนุน}

\hspace*{1.3em} %ย่อหน้า
Web Application คือการพัฒนาระบบงานบนเว็บ ข้อมูลต่าง ๆ ในระบบสามารถเข้าถึงได้ทั้งเครือข่ายภายในหรือเครือข่ายสาธารณะ ซึ่งทำให้ใช้งานง่ายผ่านเว็บบราวเซอร์ โดยในการใช้งานนั้นไม่จำเป็นต้องติดตั้งโปรแกรมใด ๆ เพิ่มเติม ทำให้เหมาะสำหรับการใช้งานในสภาพแวดล้อมที่มีผู้ใช้งานหลากหลาย

\hspace*{1.3em} %ย่อหน้า
\textbf{คุณสมบัติสำคัญของ Web Application สำหรับการจัดการระบบคลัสเตอร์}
\begin{mycustomitem}
    \item \textbf{Cross-platform Compatibility} สามารถใช้งานได้บนระบบปฏิบัติการต่าง ๆ โดยไม่ต้องพัฒนาแยกสำหรับแต่ละแพลตฟอร์ม
    \item \textbf{Centralized Management} ช่วยให้สามารถจัดการระบบจากจุดเดียว ลดความซับซ้อนในการดูแลระบบหลายเครื่อง
    \item \textbf{Role-Based Access Control (RBAC)} สามารถกำหนดสิทธิ์การเข้าถึงตามบทบาทของผู้ใช้ เช่น นักศึกษา อาจารย์ และผู้ดูแลระบบ
    \item \textbf{Real-time Monitoring} สามารถแสดงสถานะของระบบแบบเรียลไทม์ ช่วยในการติดตามและแก้ไขปัญหาได้รวดเร็ว
\end{mycustomitem}

\hspace*{1.3em} %ย่อหน้า
\textbf{Cloud-Init และการตั้งค่าเครื่องเสมือนอัตโนมัติ}

\hspace*{1.3em} %ย่อหน้า
Cloud-Init คือการตั้งค่าที่ถูกเตรียมไว้ล่วงหน้าสำหรับการจัดเตรียมเครื่องจำลองใหม่โดยอัตโนมัติ ซึ่งรวมถึงการสร้างข้อมูลประจำตัวผู้ใช้และการตั้งค่าระบบขั้นพื้นฐาน ด้วยการสร้างไฟล์ cloud-config ที่ปรับให้เหมาะกับเครื่องเสมือนแต่ละเครื่องอย่างไดนามิก

\begin{mycustomitem}
    \item \textbf{การตั้งค่าผู้ใช้งานอัตโนมัติ} สามารถสร้างบัญชีผู้ใช้ กำหนดรหัสผ่าน และตั้งค่า SSH keys ได้โดยอัตโนมัติ
    \item \textbf{การติดตั้งซอฟต์แวร์เริ่มต้น} สามารถกำหนดให้ติดตั้งแพ็กเกจต่าง ๆ ที่จำเป็นสำหรับการใช้งานเบื้องต้น
    \item \textbf{การตั้งค่าเครือข่าย} สามารถกำหนดค่า IP Address, DNS, และการตั้งค่าเครือข่ายอื่น ๆ ได้
    \item \textbf{การรันสคริปต์เริ่มต้น} สามารถกำหนดให้รันคำสั่งหรือสคริปต์ต่าง ๆ หลังจากที่ระบบบูตขึ้นมาแล้ว
\end{mycustomitem}

\hspace*{1.3em} %ย่อหน้า
\textbf{Reverse Proxy และการจัดการการเข้าถึง}

\hspace*{1.3em} %ย่อหน้า
Reverse Proxy คือตัวกลางที่ส่งข้อมูลเครือข่ายที่แลกเปลี่ยนข้อมูลระหว่างไคลเอนต์และเซิร์ฟเวอร์ โดยจะรับคำขอจากไคลเอนต์บนเครือข่ายสาธารณะ (อินเทอร์เน็ต) ส่งต่อคำขอไปยังเซิร์ฟเวอร์ในเครือข่ายภายใน จากนั้นจึงส่งคำตอบที่สร้างโดยเซิร์ฟเวอร์ไปยังไคลเอนต์

\begin{mycustomitem}
    \item \textbf{Port Forwarding} ช่วยให้สามารถเข้าถึงบริการต่าง ๆ ในเครื่องเสมือนจากภายนอกได้โดยผ่านพอร์ตที่กำหนด
    \item \textbf{Load Balancing} สามารถกระจายโหลดไปยังเซิร์ฟเวอร์หลายตัวเพื่อเพิ่มประสิทธิภาพ
    \item \textbf{SSL Termination} สามารถจัดการการเข้ารหัส SSL/TLS ได้ที่ตัว Proxy โดยไม่ต้องตั้งค่าในแต่ละเซิร์ฟเวอร์
    \item \textbf{Security} ช่วยป้องกันการเข้าถึงโดยตรงไปยังเซิร์ฟเวอร์ภายใน และสามารถกำหนดกฎการเข้าถึงได้
\end{mycustomitem}

\hspace*{1.3em} %ย่อหน้า
สรุปได้ว่า การนำเทคโนโลยี Virtualization, Cloud Computing, Container Orchestration, Web Application และเทคโนโลยีสนับสนุนต่าง ๆ มาใช้ร่วมกันในการพัฒนาแพลตฟอร์มคลัสเตอร์สำหรับการศึกษา จะช่วยให้สามารถสร้างสภาพแวดล้อมการเรียนการสอนที่มีประสิทธิภาพ ยืดหยุ่น และตอบสนองความต้องการของผู้ใช้งานได้หลากหลายรูปแบบ โดยเฉพาะอย่างยิ่งในสาขาเทคโนโลยีสารสนเทศที่ต้องการการฝึกปฏิบัติในสภาพแวดล้อมที่เสมือนจริง


%*******************************************
% งานวิจัยที่เกี่ยวข้อง
%*******************************************
\section{รายงานการค้นคว้า การศึกษา หรืองานวิจัยที่เกี่ยวข้อง}

\subsection{การพัฒนาระบบ Virtualization สำหรับการศึกษา}
\hspace*{1.3em} %ย่อหน้า
การศึกษาวิจัยเรื่อง \textbf{"Guide to Security for Full Virtualization Technologies"} โดย Karen Scarfone, Murugiah Souppaya, และ Paul Hoffman ได้นำเสนอแนวทางและมาตรฐานด้านความปลอดภัยสำหรับเทคโนโลยี Virtualization แบบเต็มรูปแบบ \parencite{GuideToSecurityforFullVT}. งานวิจัยนี้มีความสำคัญต่อการพัฒนาแพลตฟอร์มคลัสเตอร์เนื่องจากให้ข้อมูลเกี่ยวกับการจำแนกประเภทของ Virtualization ตามชั้นสถาปัตยกรรมคอมพิวเตอร์ และแนวทางการรักษาความปลอดภัยที่จำเป็น ซึ่งเป็นองค์ความรู้พื้นฐานสำคัญในการออกแบบระบบที่ปลอดภัยและเหมาะสมสำหรับสภาพแวดล้อมการศึกษา

\subsection{การประยุกต์ใช้ Cloud Computing ในการศึกษา}
\hspace*{1.3em} %ย่อหน้า
Danairat T. ได้นำเสนอเอกสารเรื่อง \textbf{"Cloud Computing Technology The Architecture Overview"} ซึ่งให้ภาพรวมของสถาปัตยกรรม Cloud Computing และแนวคิดการประมวลผลแบบคลาวด์ที่ช่วยให้สามารถเข้าถึงทรัพยากรการประมวลผลที่กำหนดค่าได้ตามต้องการ \parencite{CloudComputingTechnologyTheArch}. งานศึกษานี้เป็นแนวทางสำคัญในการทำความเข้าใจรูปแบบการให้บริการ เช่น IaaS, PaaS, และ SaaS ซึ่งเป็นแนวคิดที่นำมาประยุกต์ใช้ในการพัฒนาแพลตฟอร์มคลัสเตอร์เพื่อให้บริการเครื่องเสมือนแก่ผู้ใช้งานในสถาบันการศึกษา

\subsection{การพัฒนาระบบติดตามความก้าวหน้าสำหรับเอกสารโครงงานพิเศษ}
\hspace*{1.3em} %ย่อหน้า
Rawiporn Suamsiri และ Kamonwan Janmanee ได้พัฒนาระบบ \textbf{"Progress Tracking System for Special Project Documents"} ซึ่งเป็นระบบ Web Application ที่ช่วยในการติดตามความก้าวหน้าของการจัดทำเอกสารโครงงานพิเศษ \parencite{suamsiri2024_progress_tracking}. งานวิจัยนี้แสดงให้เห็นถึงแนวทางการพัฒนา Web Application ที่ใช้งานง่ายผ่านเว็บบราวเซอร์โดยไม่จำเป็นต้องติดตั้งโปรแกรมเพิ่มเติม ซึ่งเป็นแนวคิดที่สอดคล้องกับการพัฒนาส่วนติดต่อผู้ใช้งานสำหรับแพลตฟอร์มคลัสเตอร์ในโครงงานนี้

\subsection{การใช้งาน Container Technologies ในสถาปัตยกรรมต่าง ๆ}
\hspace*{1.3em} %ย่อหน้า
งานวิจัยของ SHAHIDULLAH KAISER, MD.SADUNHAQ, ALI AMAN TOSUN, และ TURGAY KORKMAZ เรื่อง \textbf{"Container Technologies for ARM Architecture: A Comprehensive Survey of the State-of-the-Art"} นำเสนอการสำรวจเทคโนโลยี Container บนสถาปัตยกรรม ARM อย่างครอบคลุม \parencite{kaiser2022_container_arm}. แสดงให้เห็นถึงการพัฒนาและการใช้งาน Container ที่มีประสิทธิภาพสูง ใช้ทรัพยากรน้อย และสามารถดำเนินการไมโครเซอร์วิสได้อย่างมีประสิทธิภาพ ซึ่งเป็นข้อมูลสำคัญในการออกแบบระบบที่ใช้ Container เป็นส่วนหนึ่งของโครงสร้างพื้นฐาน

\subsection{Container Orchestration ในอุตสาหกรรมยานยนต์}
\hspace*{1.3em} %ย่อหน้า
Naresh Nayak, Dennis Grewe, และ Sebastian Schildt ได้นำเสนอการศึกษาเรื่อง \textbf{"Automotive Container Orchestration: Requirements, Challenges and Open Directions"} ซึ่งอธิบายแนวคิดของ Container Orchestration ในการปรับใช้และจัดการแอปพลิเคชันที่เป็น Container ในคลัสเตอร์ \parencite{nayak2023_automotive_orchestration}. งานวิจัยนี้ให้ความรู้เกี่ยวกับการทำงานของโหนดมาสเตอร์และโหนดเวิร์กเกอร์ รวมถึงคุณสมบัติสำคัญ เช่น High Availability, Auto Scaling, และ Load Balancing ซึ่งเป็นแนวคิดที่นำมาประยุกต์ใช้ในการพัฒนาระบบจัดการคลัสเตอร์ในโครงงานนี้

\subsection{การสร้างโครงสร้างพื้นฐานคลาวด์สำหรับการจัดตารางเครื่องเสมือน}
\hspace*{1.3em} %ย่อหน้า
Jonathan Chya และ Xunfei Jiang ได้นำเสนอการศึกษาเรื่อง \textbf{"Building a Cloud Infrastructure for Virtual Machine Scheduling in Datacenters"} ซึ่งเป็นการศึกษาเกี่ยวกับการสร้างโครงสร้างพื้นฐานคลาวด์สำหรับการจัดตารางเครื่องเสมือนในศูนย์ข้อมูล รวมถึงการใช้ Cloud-Init ในการจัดเตรียมเครื่องเสมือนใหม่โดยอัตโนมัติ \parencite{chya2024_vm_scheduling}. งานวิจัยนี้แสดงให้เห็นถึงการใช้ไฟล์ cloud-config ที่ปรับให้เหมาะกับเครื่องเสมือนแต่ละเครื่องอย่างไดนามิก ซึ่งช่วยเพิ่มประสิทธิภาพด้านพลังงาน ความเร็ว และการปรับแต่งกระบวนการจัดเตรียมเครื่องเสมือนได้อย่างมาก

\subsection{การพัฒนาสถาปัตยกรรม Reverse Proxy}
\hspace*{1.3em} %ย่อหน้า
Yu TAO และ Gang CHEN ได้เผยแพร่งานวิจัยเรื่อง \textbf{"An Extensible Universal Reverse Proxy Architecture"} ซึ่งนำเสนอสถาปัตยกรรม Reverse Proxy ที่มีความยืดหยุ่นและสามารถขยายได้ \parencite{tao2016_reverse_proxy}. งานวิจัยนี้อธิบายการทำงานของ Reverse Proxy ในการรับคำขอจากไคลเอนต์บนเครือข่ายสาธารณะและส่งต่อไปยังเซิร์ฟเวอร์ในเครือข่ายภายใน รวมถึงความสามารถในการแคชทรัพยากรเพื่อเพิ่มประสิทธิภาพ ซึ่งเป็นแนวคิดสำคัญที่นำมาใช้ในการพัฒนาระบบ Port Forwarding และการจัดการการเข้าถึงเครื่องเสมือนจากภายนอกในโครงงานนี้
